Mô phỏng Monte Carlo dựa trên Copula là một công cụ mạnh mẽ trong quản trị rủi ro định lượng, giúp hiểu sâu hơn về cấu trúc phụ thuộc của các biến ngẫu nhiên. Kỹ thuật xây dựng mô hình đi từ dưới lên này cho phép các nhà quản lý rủi ro linh hoạt kết hợp cấu trúc phân phối của các biến biên độc lập với các mô hình Copula đặc trưng nhằm mô phỏng chính xác hành vi rủi ro toàn danh mục. Việc tạo lập các phân phối liên kết Meta như Meta-Gaussian hay Meta-t đóng vai trò cốt lõi phục vụ phân tích thử nghiệm và đánh giá mức độ nhạy cảm của hệ thống trước sự biến đổi từ các tham số phụ thuộc đuôi. Trong quản trị rủi ro tín dụng, mô phỏng Copula được ứng dụng để xác định phân phối tổn thất của danh mục vay hoặc trái phiếu thông qua việc mô phỏng thời gian vỡ nợ đồng thời của các đối tác, đây là phương pháp tiêu chuẩn để định giá các sản phẩm phái sinh tín dụng phức tạp như basket derivatives và các tầng của CDO (Collateralized Debt Obligations). 

Hơn nữa, việc mô phỏng các cấu trúc này cho phép các nhà giao dịch tính toán tương quan tài sản ngầm định từ giá thị trường của các tầng CDO, giúp so sánh rủi ro giữa các điểm đính kèm và nhóm tài sản khác nhau. Để cải thiện hiệu suất khi mô phỏng các biến cố cực đoan như tổn thất lớn trong các tầng CDO cao, kỹ thuật lấy mẫu ưu tiên thường được kết hợp với mô phỏng Copula để giảm biến số và tăng độ chính xác của ước lượng. Đối với các bài toán ước tính rủi ro đa kỳ hạn, mô phỏng Monte Carlo từ Copula cho phép xác định các thước đo rủi ro như VaR và expected shortfall cho các kỳ hạn dài bằng các mô phỏng các đường dẫn thực tế của các yếu tố rủi ro thay vì dựa vào các quy tắc tỷ lệ đơn giản. Ngoài ra, trong các mô hình actuarial của rủi ro hoạt động và bảo hiểm, mô phỏng được sử dụng để tính toán tổng mức tổn thất bằng cách kết hợp sự phụ thuộc giữa tần suất xuất hiện và mức độ nghiêm trọng của các loại tổn thất khác nhau.